\chapter{The Algorithms}
\label{ch:algorithms}

Everything so far defined \emph{what} to compute. This chapter gives the
numerical engines: iterative eigensolvers for the linear problem at fixed
potential, and mixing schemes for the nonlinear self-consistency, with full
pseudocode. Notation: at each $\kk$, $H_{\kk}$ is the (matrix-free)
Hamiltonian on spinor coefficient vectors of length $2N_G$; we seek the
lowest $N_b$ eigenpairs.

\section{Setting up a $\kk$-point: the basis}
\label{sec:setup}

Everything at a given $\kk$ is expressed on its own plane-wave sphere.
Building it is pure integer bookkeeping, done once per geometry.

\begin{algorithm}[H]
\caption{Plane-wave basis and projectors at one $\kk$-point}
\label{alg:basis}
\begin{algorithmic}[1]
\Function{BuildBasis}{$\kk_{\mathrm{frac}}$, $E_{\mathrm{cut}}$}
\State $M_i \gets \big\lceil \sqrt{2E_{\mathrm{cut}}}\,|\mathbf a_i|/2\pi
\big\rceil + 1$ \Comment{no $\GG$ in the sphere has $|m_i| > M_i$}
\State \textbf{for all} integer triples $\mathbf m$ with $|m_i|\le M_i$:
$\qq \gets (\kk_{\mathrm{frac}} + \mathbf m)\,B$;\quad keep $\mathbf m$ if
$\tfrac12|\qq|^2 \le E_{\mathrm{cut}}$
\State store for each kept plane wave: $\mathbf m$;\ $\qq$;\ $T_{\qq} =
\tfrac12|\qq|^2$;\ FFT-box index
$\mathrm{idx} = \mathrm{ravel}(m_1 \bmod N_1,\, m_2 \bmod N_2,\, m_3 \bmod N_3)$
\State projectors $\tilde\beta^{a}_{i\ell m}(\qq)$ from
Eqs.~\eqref{eq:betaG}, \eqref{eq:hghprojG}, Table~\ref{tab:realylm}
\Comment{array $[N_{\mathrm{proj}}, N_G]$}
\State \Return basis record (padded to $N_G^{\max}$ with a mask when
batching over $\kk$)
\EndFunction
\end{algorithmic}
\end{algorithm}

Here $B$ is the matrix whose rows are $\mathbf b_i$. The ``$\bmod N_i$''
places negative frequencies where the FFT library expects them (wrapped
order). Sort the kept plane waves in any fixed order you like; nothing
depends on it. Two checks: the count should be close to
$\Omega(2E_{\mathrm{cut}})^{3/2}/6\pi^2$, and at $\Gamma$ the set must be
symmetric under $\mathbf m\to-\mathbf m$.

\section{Applying the Hamiltonian}
\label{sec:applyh}

The eigensolver below never sees a matrix, only this function. It is
Chapter~\ref{ch:hamiltonian} condensed into one routine; $C$ holds any
number of spinor vectors (bands) at one $\kk$.

\begin{algorithm}[H]
\caption{Matrix-free Hamiltonian application at one $\kk$-point}
\label{alg:applyh}
\begin{algorithmic}[1]
\Function{ApplyH}{$C[n, s, \GG]$;\ basis;\ $V_{\mathrm{eff,loc}}(\rr)$,
$\mathbf B_{\mathrm{xc}}(\rr)$ on the grid}
\State $Y \gets T_{\qq}\, C$ \Comment{kinetic: elementwise, same for both $s$}
\State $u[n,s,\rr] \gets \mathrm{IFFT}\big(\text{scatter } C[n,s,\cdot]
\text{ into a zero box at idx}\big)$ \Comment{to real space}
\State $v[n,s,\rr] \gets V_{\mathrm{eff,loc}}(\rr)\,u[n,s,\rr]
+ \sum_{s'}\big(\mathbf B_{\mathrm{xc}}(\rr)\cdot\bm\sigma\big)_{ss'}\,
u[n,s',\rr]$ \Comment{second term absent for gold}
\State $Y \mathrel{+}= \text{gather } \mathrm{FFT}(v) \text{ at idx}$
\Comment{back to the sphere; discard everything outside it}
\State $P[n,s,p] \gets \sum_{\GG} \tilde\beta_p^{*}(\GG)\, C[n,s,\GG]$
\Comment{projections, $p = (a,i,\ell,m)$}
\State $Q[n,s,p] \gets \sum_{s'p'} D_{(p s),(p' s')}\, P[n,s',p']$
\Comment{coupling matrix, Section~\ref{sec:lsmatrix}}
\State $Y \mathrel{+}= \sum_p \tilde\beta_p(\GG)\, Q[n,s,p]$
\State \Return $Y$
\EndFunction
\end{algorithmic}
\end{algorithm}

The FFT normalizations cancel between lines 3 and 5 (use the library's
unnormalized forward and $1/N$-normalized inverse transform and forget
about them here). Test it before anything else: for random vectors
$\phi,\psi$ the two numbers $\braket{\phi|\hat H\psi}$ and
$\braket{\psi|\hat H\phi}^*$ must agree to rounding error, and for a small
cutoff the matrix assembled by applying \textsc{ApplyH} to unit vectors
must have the same eigenvalues as the iterative solver returns.

\section{Why iterative eigensolvers}

$H_{\kk}$ has dimension $2N_G \sim 5\times10^3$; dense diagonalization
($O(N^3) \sim 10^{11}$ flops \emph{per $\kk$-point per SCF step}) is possible
but wasteful, and becomes impossible for larger cells you will want later. We
need only the lowest $N_b \sim 24$ of $\sim 5\times10^3$ eigenpairs of a Hermitian
operator we can apply fast. That is the textbook setting for
\emph{preconditioned block Krylov / subspace iteration} methods. Two
GPU-relevant choices:

\begin{itemize}[leftmargin=2em]
\item \textbf{Blocked Davidson}: expands a subspace with preconditioned
residuals; robust, standard in plane-wave codes.
\item \textbf{LOBPCG} (Locally Optimal Block Preconditioned Conjugate
Gradient): keeps a fixed-size subspace $[X, W, P]$ (current block,
preconditioned residuals, previous search directions) and Rayleigh--Ritz
minimizes within it. Fixed memory, fixed shapes, everything expressed as
dense GEMMs on tall matrices --- \emph{the} eigensolver for JAX/GPU. We make
it the primary recommendation.
\end{itemize}

\section{The preconditioner}

Iterative methods converge at a rate set by the spread of the spectrum. A
plane-wave Hamiltonian's spectrum is dominated by kinetic energy
$\tfrac12|\kk+\GG|^2$, which ranges up to $E_{\mathrm{cut}}$; but we only
care about the lowest few eV. The cure is a diagonal preconditioner that
approximately inverts the kinetic term at high $\GG$ while leaving low $\GG$
alone. The classic Teter--Payne--Allan form, applied elementwise to a
residual $r(\GG)$ with $x = \tfrac12|\kk+\GG|^2 / E_{\mathrm{kin}}(\psi)$:
\begin{equation}
K(x) = \frac{27 + 18x + 12x^2 + 8x^3}
{27 + 18x + 12x^2 + 8x^3 + 16 x^4},
\label{eq:teter}
\end{equation}
smoothly interpolating between $1$ (small $x$) and $1/(2x)$ (large $x$);
$E_{\mathrm{kin}}(\psi) = \sum_{\GG s}\tfrac12|\kk{+}\GG|^2 |c^s(\GG)|^2$ is
the kinetic energy of the band being corrected.
A simpler alternative, $K = 1/\max\big(1, \tfrac12|\kk+\GG|^2 - \tau\big)$
with shift $\tau \approx$ current Ritz value, also works. The preconditioner
is elementwise and free; it typically cuts iteration counts by
$5$--$10\times$.

\section{LOBPCG}

The idea in one paragraph. Suppose $X$ holds your current guesses for the
lowest $N_b$ eigenvectors. The residual $R = HX - X\Lambda$ tells you in
which direction each guess is wrong. Rather than taking a step of some
guessed length along $R$, build the small subspace spanned by the current
vectors $X$, the (preconditioned) residuals $W$, and the previous step's
update directions $P$, and solve the eigenproblem \emph{exactly inside that
subspace} --- a $3N_b\times3N_b$ dense diagonalization, called the
\textbf{Rayleigh--Ritz} step. The best $N_b$ vectors of the subspace
become the new $X$. Because the subspace contains the old $X$, the Ritz
values can only go down; because it contains $P$, the method has the
memory that makes conjugate gradients fast.

Written naively, with $S = [X, W, P]$ taken as it comes, the Rayleigh--Ritz
step is a \emph{generalized} eigenproblem $S^\dagger H S\, y = \theta\,
S^\dagger S\, y$, and it is the weak point of the method: near convergence
$W$ and $P$ shrink and become almost parallel to each other, the overlap
matrix $S^\dagger S$ becomes numerically singular, and the iteration falls
apart (details below). The cure is to keep the three blocks
\emph{orthonormal and mutually orthogonal at all times}, so that
$S^\dagger S = \mathbf 1$ and the Rayleigh--Ritz step is an ordinary
Hermitian eigenproblem. One small utility does all the orthonormalizing.

\begin{algorithm}[H]
\caption{Orthonormalize the columns of $V$, discarding (zeroing) linearly
dependent directions; optionally carry $HV$ along}
\label{alg:ortho}
\begin{algorithmic}[1]
\Function{Orthonormalize}{$V$, $\tau$ [, $HV$]}
\State $d_j \gets \|V_{:,j}\|$ (set $d_j = 1$ where the column is zero)
\Comment{column scales}
\State $G \gets \mathrm{diag}(d)^{-1}\, V^\dagger V\, \mathrm{diag}(d)^{-1}$
\Comment{small Hermitian matrix with unit diagonal; double precision}
\State $(\sigma, U) \gets \mathrm{eigh}(G)$
\Comment{$\sigma_j \ge 0$ measures how independent direction $j$ is}
\State $T \gets \mathrm{diag}(d)^{-1}\; U\;
\mathrm{diag}\big(\sigma_j > \tau\ ?\ \sigma_j^{-1/2} : 0\big)$
\State \Return $V\,T$ [and $HV\,T$]
\Comment{columns are orthonormal, or exactly zero}
\EndFunction
\end{algorithmic}
\end{algorithm}

Because the same matrix $T$ multiplies $V$ and $HV$, the pair stays
consistent ($HV$ remains the image of $V$ under $\hat H$) without a new
Hamiltonian application. Dependent directions are not removed --- that
would change array shapes --- but replaced by zero columns, which the
eigensolver then ignores. There is no failure branch, unlike a Cholesky
factorization, which simply breaks down on a singular matrix.

\begin{algorithm}[H]
\caption{LOBPCG for the lowest $N_b$ spinor bands at one $\kk$
(all $\kk$ run batched)}
\label{alg:lobpcg}
\begin{algorithmic}[1]
\State $X \gets$ initial block $[2N_G \times N_b]$ (previous SCF
step's eigenvectors; first step: random complex numbers divided by
$1 + \tfrac12|\kk{+}\GG|^2$)
\State $X \gets$ \Call{Orthonormalize}{$X$, $10^{-12}$}, twice;\quad
$HX \gets$ \Call{ApplyH}{$X$}
\State $(\Lambda, Y) \gets \mathrm{eigh}(X^\dagger HX)$;\quad
$X \gets XY$;\ $HX \gets HX\,Y$ \Comment{start from Ritz vectors}
\State $P \gets 0$;\ $HP \gets 0$ \Comment{same shape as $X$}
\For{$t = 1, \dots, t_{\max}$}
\State $R \gets HX - X\,\mathrm{diag}(\Lambda)$;\quad
$r_n \gets \|R_{:,n}\|$ \Comment{residual of each band}
\State \textbf{if} $r_n < \epsilon$ for all bands that are needed
\textbf{then stop}
\State $W \gets K \circ R$ \Comment{preconditioner Eq.~\eqref{eq:teter},
each column with its own $E_{\mathrm{kin}}$}
\State $W_{:,n} \gets 0$ where $r_n < \epsilon/10$
\Comment{``soft locking'': converged bands get no new direction}
\State \textbf{twice:}\ \ $W \gets W - Q\,(Q^\dagger W)$ with $Q = [X, P]$;\quad
$W \gets$ \Call{Orthonormalize}{$W$, $10^{-10}$}
\State $HW \gets$ \Call{ApplyH}{$W$} \Comment{the only Hamiltonian
application per iteration}
\State $S \gets [X, W, P]$;\quad $A \gets S^\dagger\,[HX, HW, HP]$;\quad
$A \gets \tfrac12(A + A^\dagger)$ \Comment{$3N_b\times3N_b$, double precision}
\State $A_{jj} \gets A_{jj} + \Lambda_{\mathrm{big}}$ for every $j$ whose
column of $S$ is zero \Comment{$\Lambda_{\mathrm{big}} = 10^6$\,Ha: park
empty directions out of the way}
\State $(\Theta, Y) \gets \mathrm{eigh}(A)$;\quad
$\Lambda \gets \Theta_{1..N_b}$;\quad keep the first $N_b$ columns of $Y$,
split by rows as $(Y_X; Y_W; Y_P)$
\State $P \gets W Y_W + P\, Y_P$;\quad $HP \gets HW\, Y_W + HP\, Y_P$
\Comment{the step just taken}
\State $X \gets X Y_X + P$;\quad $HX \gets HX\, Y_X + HP$
\Comment{new Ritz vectors; no extra $H$ applications}
\State \textbf{twice:}\ \ $C \gets X^\dagger P$;\ \ $P \gets P - XC$;\ \
$HP \gets HP - HX\,C$;\ \
$(P, HP) \gets$ \Call{Orthonormalize}{$P$, $10^{-6}$, $HP$}
\Comment{$P \perp$ new $X$, orthonormal}
\EndFor
\State \Return $\Lambda$ (eigenvalues, ascending), $X$ (eigenvectors)
\end{algorithmic}
\end{algorithm}

Note the bookkeeping: $HX$ and $HP$ are updated by the same linear
combinations as $X$ and $P$, so the expensive operation is applied only to
$W$ --- and it is applied \emph{after} $W$ has been orthogonalized, so that
$HW$ is exact for the vectors actually used. The first iteration needs no
special case: $P = 0$ contributes zero columns, which line 13 parks. The
new $X$ is orthonormal because $S$ is and the columns of $Y$ are.

Two remarks on the tolerance $\epsilon$. In a fixed-potential calculation
(Section~\ref{sec:nscf}) it is the accuracy you want, typically
$10^{-8}$\,Ha, and the loop stops when it is reached. Inside the SCF loop
only a fixed, small number of iterations is done per step; there set
$\epsilon = 0$, so that nothing is ever locked and there is no stopping
test --- which is also what a jitted program with a static iteration count
wants.

\paragraph{Why ``twice'' appears twice.} Lines 10 and 17 each repeat a
projection and an orthonormalization. This is not belt and braces.
\textsc{Orthonormalize} divides by $\sqrt{\sigma_j}$, which amplifies
rounding errors by up to $\tau^{-1/2}$ ($10^{5}$ for $W$, $10^{3}$ for
$P$, on top of the rescaling by the column norms, which for a converged
band are of order $10^{-10}$). After a single pass a nearly dependent
block is therefore orthonormal only to about $10^{-10}$, not $10^{-16}$.
The Rayleigh--Ritz step \emph{assumes} $S^\dagger S = \mathbf 1$, so the
defect is inherited by the new $X$ and grows from one iteration to the
next. The second pass starts from an almost orthonormal block, amplifies
nothing, and brings the defect down to rounding level. For $P$ the second
pass costs two small matrix products and no Hamiltonian application,
because $HP$ is carried along. Keep both passes: with a single pass on $P$
the iteration is fine while it converges and then, some tens of iterations
\emph{after} reaching rounding-level residuals, loses orthogonality and
diverges. (This was found the hard way: an implementation written from an
earlier version of this text diverged on the cosine-potential test of
Chapter~\ref{ch:verify}, Stage~1, while passing all tests on gold.)

\paragraph{Numerical stability --- read this before simplifying anything.}
LOBPCG as usually presented in three lines (take $S = [X, W, P]$ as is,
Cholesky-factorize $S^\dagger S$, solve) is a trap. It behaves perfectly
for the first ten or twenty iterations, which is all an SCF step uses; then,
once some bands are converged to rounding error, it \emph{diverges} --- not
to a slightly wrong answer but to eigenvalues of $-10^{3}$\,Ha and far
beyond --- usually without any error being raised. For the gold
Hamiltonian of this document, in double precision, this happens within
30--100 iterations at ordinary $\kk$-points, i.e.\ exactly in the tightly
converged band-structure and topology calculations. The mechanism is worth
understanding, because it defeats the obvious repairs too (for instance,
replacing Cholesky by an eigendecomposition of $S^\dagger S$ with a small
cutoff): when $W$ and $P$ are nearly parallel, the Ritz coefficients $Y$
become huge ($10^5$) with opposite signs, the new $P$ is a small difference
of large vectors, and the stored $HP$ --- updated by recurrence, never
recomputed --- stops being the image of $P$ under $\hat H$. That
inconsistency is amplified at every iteration, leaks into $HX$, and from
then on the ``Rayleigh quotients'' are no longer bounded below by the true
lowest eigenvalue. The design of Algorithm~\ref{alg:lobpcg} removes the
cause: all coefficient matrices are $O(1)$ because $S$ is orthonormal; the
only block whose image is obtained by recurrence, $P$, is transformed
exclusively by well-conditioned operations (line 17 discards directions
below $10^{-6}$ instead of amplifying them); and nearly dependent
directions in $W$ are discarded \emph{before} $\hat H$ is applied. In
tests on the gold Hamiltonian --- $\Gamma$, $X$, $L$, $W$, $K$ and generic
$\kk$, with and without SOC, several random starts each --- and on the
highly degenerate nearly-free-electron problems of Stage~1, with
$\epsilon = 10^{-10}$, $10^{-6}$ and $0$, this version converged in every
case and stayed converged (eigenvalues within $2\times10^{-9}$\,Ha of dense
diagonalization, $X^\dagger X = \mathbf 1$ to $10^{-13}$) for 400
iterations. One thing it does not promise: once the residuals have
reached the rounding floor ($10^{-9}$--$10^{-10}$\,Ha for a level with
many degenerate partners) they wander there instead of decreasing, so ask
for $\epsilon \ge 10^{-9}$\,Ha and stop when you get it.

\paragraph{What to expect.} For gold with $N_b = 20$ spinor bands and a
random start, the residuals of the lowest 14 bands fall below
$10^{-9}$\,Ha in 25--45 iterations; the eigenvalues are converged to
$10^{-10}$\,Ha well before that (the eigenvalue error is quadratic in the
residual). The top few bands of the block always converge slowest ---
which is why $N_b$ should exceed the number of bands you need by a buffer
of four or more.

\paragraph{How many iterations?} Inside an SCF loop you do \emph{not} need
tight convergence: 2--4 LOBPCG iterations per SCF step suffice (after
$\sim8$ in the very first step, which starts from random vectors), because
the potential will move anyway and the residuals keep shrinking from one
SCF step to the next. Warm-starting from the previous step's $X$ is what
makes the whole calculation fast.

\paragraph{Blocked Davidson (alternative).} Expand $V \gets [V, K\circ R]$ up
to a max subspace ($\sim 3 N_b$), Rayleigh--Ritz (Algorithms~\ref{alg:lobpcg}
and \ref{alg:ortho} serve as building blocks), restart with Ritz vectors.
Same kernels, but the growing/shrinking subspace means dynamic shapes ---
awkward under JIT; if you implement it, fix the schedule (e.g.\ always
exactly two expansions per restart) so shapes are static.

\section{Occupations, density, and the SCF loop}

After diagonalization at all $\kk$: find $\mu$, set occupations, and
accumulate the density matrix on the grid.

\begin{algorithm}[H]
\caption{Fermi level and occupations}
\label{alg:fermi}
\begin{algorithmic}[1]
\Function{Occupy}{$\varepsilon_{n\kk}$, $w_\kk$, $N_e$, $\sigma$}
\State $\mu_{\mathrm{lo}} \gets \min\varepsilon - 10\sigma$;\quad
$\mu_{\mathrm{hi}} \gets \max\varepsilon + 10\sigma$
\For{$80$ iterations} \Comment{fixed count: interval shrinks by $2^{80}$}
\State $\mu \gets \tfrac12(\mu_{\mathrm{lo}}+\mu_{\mathrm{hi}})$;\quad
$N(\mu) \gets \sum_{n\kk} w_\kk\, f\big((\varepsilon_{n\kk}-\mu)/\sigma\big)$
\State \textbf{if} $N(\mu) > N_e$ \textbf{then} $\mu_{\mathrm{hi}}\gets\mu$
\textbf{else} $\mu_{\mathrm{lo}}\gets\mu$
\EndFor
\State $x_{n\kk} \gets (\varepsilon_{n\kk}-\mu)/\sigma$;\quad
$f_{n\kk} \gets f(x_{n\kk})$;\quad
$S \gets \sum_{n\kk} w_\kk\, s(x_{n\kk})$
\State \Return $\mu$, $f$, $S$ \Comment{entropy term of the free energy is
$-\sigma S$}
\EndFunction
\end{algorithmic}
\end{algorithm}

\begin{algorithm}[H]
\caption{Density matrix from the spinor bands}
\label{alg:density}
\begin{algorithmic}[1]
\State $n_{\alpha\beta}(\rr) \gets 0$
\For{each $\kk$ (in chunks), each band $n$}
\State $u^{s}(\rr) \gets \dfrac{N_{\mathrm{grid}}}{\sqrt\Omega}\,
\mathrm{IFFT}\big(\text{scatter } c^{s}_{n\kk}\big)$ for $s =
\uparrow,\downarrow$ \Comment{normalization of Eq.~\eqref{eq:fourierconv}}
\State $n_{\alpha\beta}(\rr) \mathrel{+}= w_\kk f_{n\kk}\,
u^{\alpha}(\rr)\,\big[u^{\beta}(\rr)\big]^{*}$
\EndFor
\State $n \gets n_{\uparrow\uparrow}+n_{\downarrow\downarrow}$;\quad
$m_x \gets 2\,\mathrm{Re}\,n_{\uparrow\downarrow}$;\
$m_y \gets -2\,\mathrm{Im}\,n_{\uparrow\downarrow}$;\
$m_z \gets n_{\uparrow\uparrow}-n_{\downarrow\downarrow}$
\end{algorithmic}
\end{algorithm}

The check: $\frac{\Omega}{N_{\mathrm{grid}}}\sum_{\rr} n(\rr) = N_e$ to
rounding error, every iteration. For gold $\mathbf m$ comes out zero to
rounding error \emph{only} if both members of every Kramers pair are
equally occupied \emph{and} both $\kk$ and $-\kk$ are in the mesh. On a
time-reversal-reduced mesh (one point of each $\pm\kk$ pair with doubled
weight, as in the reference cases of Chapter~\ref{ch:verify}) the doubled
weights give the correct $n$ but a spurious, nonzero $\mathbf m$
(Eq.~\eqref{eq:mpair}): there, accumulate only
$n = \sum |u^\uparrow|^2 + |u^\downarrow|^2$ and set $\mathbf m = 0$
--- which is also the right choice for a first version of the code. (The
reference calculations of Chapter~\ref{ch:verify} do exactly this.)

\subsection{Why simple iteration diverges for metals: charge sloshing}

The naive update $n_{\mathrm{in}}^{(m+1)} = n_{\mathrm{out}}^{(m)}$ diverges
whenever the dielectric response is strong. To see why, perturb the input
density by $\delta n_{\mathrm{in}}$. The potential changes by
$K_{\mathrm{Hxc}}\,\delta n_{\mathrm{in}}$ (dominated by the Hartree
kernel $4\pi/G^2$), and the electrons respond to that with
$\delta n_{\mathrm{out}} = \chi_0 K_{\mathrm{Hxc}}\,\delta n_{\mathrm{in}}$,
where $\chi_0$ is the density response of the non-interacting system. So an
error in the input is multiplied by $\chi_0 K_{\mathrm{Hxc}}$ each
iteration. In a metal $\chi_0(G\to0) \to -g(\varepsilon_F)$, a finite
negative number, so for a Fourier component of wavevector $G$ the
multiplier is $\approx -k_{TF}^2/G^2$ with $k_{TF}^2 = 4\pi
g(\varepsilon_F)$: negative (the error changes sign --- charge
\emph{sloshes} back and forth) and larger than one in magnitude whenever
$G < k_{TF}$ (the oscillation grows). Simple linear mixing
$n_{\mathrm{in}} \gets n_{\mathrm{in}} + A\,R$ multiplies the error by
$1 - A\,(1 + k_{TF}^2/G^2)$ instead, which is stable only if
$A < 2/(1 + k_{TF}^2/G_{\min}^2)$ --- tiny for a long cell.

How bad is it for us? In the gold primitive cell the smallest nonzero
reciprocal vector is $|\mathbf b_i| = 1.41$\,Bohr$^{-1}$, larger than
$k_{TF} \approx 0.9$\,Bohr$^{-1}$, so even plain mixing with
$A \approx 0.5$ converges. Measured on a small gold test case (scalar
relativistic, $E_{\mathrm{cut}} = 20$\,Ha, $4^3$ mesh, iterations to
$\|R\| < 10^{-7}$): plain mixing with $A = 0.5$ needs 24 iterations, with
$A = 0.8$ it has not converged after 60, and with $A = 1$ it does not
converge at all (the $d$ electrons make gold's true screening stronger
than the free-electron estimate of $k_{TF}$); Kerker mixing with $A = 0.8$
needs 16; Kerker--Pulay needs 10. Sloshing becomes the dominant problem the
moment you build a supercell or a slab (a ten-layer slab with vacuum has
$G_{\min} \approx 0.1$\,Bohr$^{-1}$). Implement the two standard fixes now, used
together, so the code is ready:

\paragraph{(a) Kerker preconditioning.} Damp the long-wavelength part of the
density update in Fourier space:
\begin{equation}
\delta n(\GG) \;\gets\; A\, \frac{G^2}{G^2 + k_{TF}^2}\; \delta n(\GG),
\qquad k_{TF}^2 = 4\pi\, g(\varepsilon_F) \ \ (\text{use } k_{TF} \approx
1\ \mathrm{Bohr}^{-1},\ A \approx 0.8).
\label{eq:kerker}
\end{equation}
This is precisely the inverse of the factor $1 + k_{TF}^2/G^2$ found above:
after preconditioning, every wavelength converges at the same rate.

\paragraph{(b) Pulay / Anderson mixing (DIIS).} Keep the last $M \sim 8$
pairs of (input density, residual $R^{(m)} = n_{\mathrm{out}}^{(m)} -
n_{\mathrm{in}}^{(m)}$); find the linear combination of past inputs whose
extrapolated residual is smallest:
\begin{equation}
\min_{\{\gamma\}} \Big\| \sum_m \gamma_m R^{(m)} \Big\|^2
\quad \text{s.t.} \sum_m \gamma_m = 1
\;\;\Longrightarrow\;\;
\begin{pmatrix} \mathsf{B} & \mathbf 1\\ \mathbf 1^{\mathsf T} & 0
\end{pmatrix}
\begin{pmatrix} \bm\gamma \\ \lambda\end{pmatrix}
= \begin{pmatrix}\mathbf 0\\ 1\end{pmatrix},
\qquad
\mathsf B_{mm'} = \braket{R^{(m)} | R^{(m')}},
\end{equation}
then update
\begin{equation}
n_{\mathrm{in}}^{(\mathrm{new})}
= \sum_m \gamma_m\Big[ n_{\mathrm{in}}^{(m)}
+ \mathcal K\, R^{(m)} \Big],
\label{eq:pulayupdate}
\end{equation}
with $\mathcal K$ the Kerker operator \eqref{eq:kerker}. The logic: if the
residual depended linearly on the input, the combination
$\sum_m\gamma_m n_{\mathrm{in}}^{(m)}$ would have residual
$\sum_m\gamma_m R^{(m)}$, the smallest reachable from the history; one
preconditioned step is then taken from there.

\begin{algorithm}[H]
\caption{Kerker-preconditioned Pulay mixing (all arrays in Fourier space)}
\label{alg:pulay}
\begin{algorithmic}[1]
\Function{Mix}{$n_{\mathrm{in}}(\GG)$, $R(\GG)$, history}
\State append $(n_{\mathrm{in}}, R)$ to the history; drop the oldest pair if
more than $M$ are stored
\State $\mathsf B_{mm'} \gets \mathrm{Re}\sum_{\GG\ne0}
\dfrac{4\pi}{G^2}\,\big[R^{(m)}(\GG)\big]^{*} R^{(m')}(\GG)$;\quad
$\mathsf B \gets \mathsf B/\max_{mm'}|\mathsf B_{mm'}|$
\State solve the bordered system for $\bm\gamma$ by least squares
\Comment{it becomes ill-conditioned near convergence}
\State \Return $\sum_m \gamma_m\big[n^{(m)}_{\mathrm{in}}(\GG)
+ A\,\tfrac{G^2}{G^2+k_{TF}^2}\,R^{(m)}(\GG)\big]$
\EndFunction
\end{algorithmic}
\end{algorithm}

The metric $4\pi/G^2$ weights long wavelengths, matching the physics (it
is the Hartree energy of the residual). The $\GG=0$ component needs no
special treatment: every input and output density integrates to $N_e$, so
$R(\GG{=}0)=0$ identically, and because $\sum_m\gamma_m = 1$ the mixed
density keeps the correct electron count. The mixed density can dip
slightly negative at isolated grid points; the clamp inside the XC routine
(Chapter~\ref{ch:xc}) is all that is needed. With the first history entry
alone, $\gamma_1 = 1$ and the step is plain Kerker mixing. Mix all four
components $(n, \mathbf m)$ with the same $\bm\gamma$ (the $\mathbf m$
components without the Kerker factor); for gold, $\mathbf m$ stays zero but
the code path should exist.

\paragraph{Measuring convergence.} Throughout this document the residual
norm is
\begin{equation}
\|R\| = \Big[\int_\Omega \big|n_{\mathrm{out}}(\rr)-n_{\mathrm{in}}(\rr)\big|^2
\dd\rr\Big]^{1/2}
= \Big[\frac{\Omega}{N_{\mathrm{grid}}}\sum_{\rr}|R(\rr)|^2\Big]^{1/2}.
\label{eq:resnorm}
\end{equation}
Because the energy is stationary at self-consistency, its error is
\emph{quadratic} in $\|R\|$: $\|R\|\sim10^{-4}$ already gives the energy to
$\sim10^{-8}$\,Ha.

\subsection{The complete SCF driver}

\begin{algorithm}[H]
\caption{Self-consistent field for spinor plane-wave DFT}
\begin{algorithmic}[1]
\State \textbf{Geometry-time setup:} lattice, reciprocal vectors, $\GG$
spheres and masks per $\kk$, FFT grid and pad/gather indices, structure
factor, $V_{\mathrm{loc,ion}}(\GG)$, projector arrays
$\tilde\beta(\kk{+}\GG)$, $D^{\ell}$ blocks (with SOC), Ewald energy
\State $n_{\mathrm{in}}(\GG) \gets \frac{1}{\Omega}\sum_a Z_{\mathrm{ion}}\,
\ee^{-G^2 w^2/2}\,\ee^{-\ii\GG\cdot\RR_a}$ with $w = 1$\,Bohr
(a Gaussian blob per atom, integrating to $N_e$); $\mathbf m \gets 0$
\State $X_{\kk} \gets$ random initial blocks (Algorithm~\ref{alg:lobpcg},
line 1)
\For{$m = 1, \dots, m_{\max}$}
\State $\vH \gets 4\pi n_{\mathrm{in}}(\GG)/G^2$;\quad
$(\vxc, \mathbf B_{\mathrm{xc}}) \gets$ XC$(n_{\mathrm{in}}, \mathbf m, \nabla n_{\mathrm{in}})$
\State $V_{\mathrm{eff,loc}} \gets V_{\mathrm{loc,ion}} + \vH + \vxc$
(real-space grid)
\State \textbf{for all $\kk$ in parallel:} Algorithm~\ref{alg:lobpcg}
started from the current $X_\kk$, with $\epsilon = 0$ and $t_{\max} = 8$ in
the first pass, $3$ afterwards, on
$H_{\kk}[V_{\mathrm{eff,loc}}, \mathbf B_{\mathrm{xc}}]$
$\Rightarrow \varepsilon_{n\kk}, X_{\kk}$
\State $(\mu, f_{n\kk}, S) \gets$ \Call{Occupy}{$\varepsilon, w, N_e, \sigma$}
\Comment{Algorithm~\ref{alg:fermi}}
\State $(n_{\mathrm{out}}, \mathbf m_{\mathrm{out}}) \gets$ density from
$\{X_{\kk}, f\}$ \Comment{Algorithm~\ref{alg:density}}
\State $R \gets n_{\mathrm{out}} - n_{\mathrm{in}}$;\quad record
$\|R\|$ and both free energies \Comment{Eqs.~\eqref{eq:resnorm},
\eqref{eq:fdirect}, \eqref{eq:ftot}}
\State \textbf{if} $\|R\| < 10^{-7}$ \textbf{stop}
\State $n_{\mathrm{in}} \gets$ \Call{Mix}{$n_{\mathrm{in}}, R$, history}
\Comment{Algorithm~\ref{alg:pulay}, $M = 8$}
\EndFor
\State Keep $n_{\mathrm{in}}$, $V_{\mathrm{eff,loc}}$ and $\mu$; all post-processing
(band structure, DOS, parities, Wilson loops) uses
Section~\ref{sec:nscf}
\end{algorithmic}
\end{algorithm}

\paragraph{Convergence pathologies and cures.}
\begin{itemize}[leftmargin=2em]
\item \emph{Oscillating $F$, residual stuck}: smearing too small or
$\kk$-mesh too coarse; raise $\sigma$ to $0.02$\,Ha to converge, then
re-refine.
\item \emph{Slow creep}: mixing history too short, or Kerker $k_{TF}$
mismatched; try $A \in [0.3, 0.8]$.
\item \emph{Occupied/empty band flip-flops}: not enough empty bands; ensure
the highest computed band is $\gtrsim 10\sigma$ above $\mu$.
\item \emph{LOBPCG stagnation}: block too small (use $N_b \ge N_{\mathrm{occ}}
+ 8$).
\item \emph{Eigenvalues suddenly huge and negative}: the stored $HX$ or
$HP$ is no longer consistent with $X$ or $P$; you simplified
Algorithm~\ref{alg:lobpcg} (see ``Numerical stability'').
\item \emph{Electron count drifts}: the density normalization of
Algorithm~\ref{alg:density} or the $\kk$ weights are off; fix this before
looking at anything else.
\end{itemize}

\section{After self-consistency: calculations at fixed potential}
\label{sec:nscf}

Once the SCF has converged, the density and hence $V_{\mathrm{eff,loc}}$
are known and \emph{frozen}. The Kohn--Sham Hamiltonian is then a fixed
linear operator, and its eigenstates can be computed at any $\kk$
whatsoever --- not only on the SCF mesh --- at the cost of one eigensolve
each. Every remaining result of this document is obtained this way.

\begin{algorithm}[H]
\caption{Bands at arbitrary $\kk$-points in a fixed potential}
\label{alg:nscf}
\begin{algorithmic}[1]
\Function{SolveBands}{list of $\kk$;\ $V_{\mathrm{eff,loc}}$;\
$N_{\mathrm{want}}$;\ $\epsilon$}
\For{each $\kk$ (batched)}
\State basis $\gets$ \Call{BuildBasis}{$\kk$, $E_{\mathrm{cut}}$}
\Comment{same cutoff and FFT grid as the SCF}
\State run Algorithm~\ref{alg:lobpcg} from a random start with
$N_b = N_{\mathrm{want}}+6$ until the lowest $N_{\mathrm{want}}$ residuals
are below $\epsilon$
\EndFor
\State \Return $\varepsilon_{n\kk}$ and coefficients $c^s_{n\kk}(\GG)$ for
$n \le N_{\mathrm{want}}$
\EndFunction
\end{algorithmic}
\end{algorithm}

Use double precision and $\epsilon \sim 10^{-8}$\,Ha here. Its uses:
\begin{itemize}[leftmargin=2em]
\item \textbf{Band structure.} A path
$\Gamma$--$X$--$W$--$L$--$\Gamma$--$K$ through the FCC high-symmetry points
of Section~\ref{sec:bzintegration}, 20--30 points per segment (interpolate
linearly in fractional coordinates; plot against the cumulative Cartesian
path length), energies measured from $\mu$.
\item \textbf{Density of states.} On a mesh at least as fine as the SCF
mesh,
\begin{equation}
g(\varepsilon) = \sum_{n\kk} w_{\kk}\,
\delta_\sigma(\varepsilon - \varepsilon_{n\kk}),
\end{equation}
with $\delta_\sigma$ a normalized Gaussian of width $\sim 0.005$\,Ha. The
$d$-band complex appears as a massive peak $2$--$7$\,eV below $E_F$, split
by SOC.
\item \textbf{Wavefunction diagnostics.} Parities at the eight TRIM
(Section~\ref{sec:symmetry}) and overlaps between neighbouring
$\kk$-points (Chapter~\ref{ch:topology}) need the coefficients, not just
the eigenvalues --- which is why \textsc{SolveBands} returns them.
\end{itemize}
One caution: eigenvectors within a degenerate level --- and in gold every
level is at least two-fold degenerate --- are returned in an arbitrary
orthonormal combination, with arbitrary overall phases, different at every
$\kk$. Any quantity you compute from wavefunctions must be insensitive to
this freedom. Eigenvalues, densities, and traces or eigenvalues of
matrices built from \emph{complete} degenerate groups are; individual
matrix elements are not.

\begin{gpubox}{Parallel structure of the whole SCF step}
One SCF iteration is: [pointwise XC kernel] $\to$ [a few FFTs for
$\vH$, gradients] $\to$ [batched LOBPCG: per iteration one batched
$\hat H$-apply $=$ $2 N_k N_b$ FFT pairs $+$ a handful of batched GEMMs,
then one $3N_b$-dim and four $N_b$-dim \texttt{eigh} per $\kk$ (batched)] $\to$ [device-wide reductions
for $\mu$] $\to$ [$2 N_k N_b$ IFFTs $+$ fused reduction for the density]
$\to$ [Fourier-space mixing]. Every stage is either elementwise, an FFT
batch, a GEMM batch, or a small batched \texttt{eigh} --- there is no
sequential inner structure anywhere except the SCF iteration count itself
($\sim 20$--$40$ iterations) and the handful of LOBPCG steps inside each.
\end{gpubox}
